% ============================================================================
% APENDICES
% ============================================================================

\appendix

\section{Detalles de implementaci\'on del pipeline}
\label{appendix:pipeline}

\subsection{Arquitectura del software}

El pipeline se implement\'o como un paquete Python instalable (\texttt{alphafold-comparison})
con la siguiente estructura modular:

\begin{lstlisting}[language=bash,caption={Estructura del paquete Python}]
src/alphafold_comparison/
    __init__.py
    __main__.py          # CLI dispatcher
    config.py            # Configuracion centralizada (.env)
    download/
        downloader.py    # Descarga PDB + AlphaFold
    preprocessing/
        processor.py     # Mapeo SIFTS + superposicion Kabsch
    analysis/
        structural.py    # Analisis RMSD estadistico
        atomic.py        # Analisis por elemento quimico (CHONSP)
    visualization/
        plots.py         # Funciones de visualizacion
\end{lstlisting}

\subsection{Dependencias principales}

\begin{table}[H]
\centering
\caption{Dependencias del proyecto y sus versiones m\'inimas.}
\label{tab:dependencies}
\begin{tabular}{@{}lll@{}}
\toprule
\textbf{Paquete} & \textbf{Versi\'on} & \textbf{Uso} \\
\midrule
NumPy & $\geq$ 1.24.0 & C\'alculo num\'erico y \'algebra lineal \\
Pandas & $\geq$ 2.0.0 & Manipulaci\'on de datos tabulares \\
SciPy & $\geq$ 1.11.0 & Tests estad\'isticos \\
BioPython & $\geq$ 1.81 & Parsing PDB, superposici\'on Kabsch \\
Matplotlib & $\geq$ 3.7.0 & Visualizaci\'on est\'atica \\
Seaborn & $\geq$ 0.12.0 & Visualizaci\'on estad\'istica \\
Requests & $\geq$ 2.31.0 & Descargas HTTP \\
tqdm & $\geq$ 4.65.0 & Barras de progreso \\
python-dotenv & $\geq$ 1.0.0 & Variables de entorno \\
\bottomrule
\end{tabular}
\end{table}

\subsection{Par\'ametros de configuraci\'on}

\begin{table}[H]
\centering
\caption{Par\'ametros configurables del pipeline.}
\label{tab:params}
\begin{tabular}{@{}llp{7cm}@{}}
\toprule
\textbf{Par\'ametro} & \textbf{Valor} & \textbf{Descripci\'on} \\
\midrule
\texttt{DEFAULT\_WORKERS} & 8 & Procesos paralelos para an\'alisis \\
\texttt{DOWNLOAD\_WORKERS} & 15 & Hilos para descargas HTTP \\
\texttt{MIN\_RESIDUES} & 3 & M\'inimo de residuos comunes para superposici\'on \\
\texttt{MIN\_ATOMS\_PER\_RES} & 1 & M\'inimo de \'atomos por residuo \\
\texttt{MIN\_ALIGNMENT\_ATOMS} & 3 & M\'inimo de \'atomos para alineamiento \\
\texttt{MIN\_COMMON\_RESIDUES} & 10 & M\'inimo de residuos comunes para an\'alisis at\'omico \\
\texttt{MIN\_DIVERSE\_UNIPROTS} & 10.000 & Prote\'inas \'unicas objetivo en dataset diverso \\
\bottomrule
\end{tabular}
\end{table}


\section{Mapeo SIFTS: detalles t\'ecnicos}
\label{appendix:sifts}

\subsection{Estructura del archivo SIFTS}

El archivo \texttt{pdb\_chain\_uniprot.csv} de SIFTS contiene las siguientes columnas
relevantes para el mapeo:

\begin{table}[H]
\centering
\caption{Columnas del archivo SIFTS utilizadas en el pipeline.}
\label{tab:sifts_columns}
\begin{tabular}{@{}llp{7cm}@{}}
\toprule
\textbf{Columna} & \textbf{Ejemplo} & \textbf{Descripci\'on} \\
\midrule
\texttt{PDB} & 3lzt & Identificador PDB (4 caracteres) \\
\texttt{CHAIN} & A & Cadena del PDB que corresponde a la prote\'ina \\
\texttt{SP\_PRIMARY} & P00698 & Identificador UniProt de la prote\'ina \\
\texttt{PDB\_BEG} & 1 & Primer residuo en numeraci\'on PDB \\
\texttt{PDB\_END} & 129 & \'Ultimo residuo en numeraci\'on PDB \\
\texttt{SP\_BEG} & 19 & Primer residuo en numeraci\'on UniProt \\
\texttt{SP\_END} & 147 & \'Ultimo residuo en numeraci\'on UniProt \\
\bottomrule
\end{tabular}
\end{table}

\subsection{Ejemplo: lisozima de clara de huevo (P00698)}

La lisozima ilustra el problema de numeraci\'on:

\begin{itemize}
    \item \textbf{UniProt P00698}: 147 residuos (posiciones 1--147). Los residuos 1--18
    corresponden al p\'eptido se\~nal, que se elimina en la prote\'ina madura.

    \item \textbf{PDB 3LZT}: Contiene la prote\'ina madura (129 residuos), numerados
    como 1--129 en el archivo PDB.

    \item \textbf{AlphaFold}: Predice la secuencia completa UniProt (1--147), pero la
    prote\'ina madura corresponde a las posiciones 19--147.

    \item \textbf{SIFTS}: Mapea PDB\_BEG=1, PDB\_END=129 $\rightarrow$ SP\_BEG=19,
    SP\_END=147. Esto permite saber que el residuo 1 del PDB corresponde al residuo
    19 de AlphaFold.

    \item \textbf{Sin mapeo}: Se compara el residuo PDB 1 con el residuo AlphaFold 1
    (que es el p\'eptido se\~nal, a $\sim$18\AA{} de distancia).

    \item \textbf{Con mapeo}: Se compara el residuo PDB 1 con el residuo AlphaFold 19
    (el residuo correcto), obteniendo RMSD $\sim$0,44\AA.
\end{itemize}


\section{Estad\'isticas complementarias}
\label{appendix:stats}

\subsection{Distribuci\'on de RMSD por rangos}

\begin{table}[H]
\centering
\caption{Distribuci\'on de RMSD con mapeo SIFTS correcto (5.115 pares validados).}
\label{tab:rmsd_distribution}
\begin{tabular}{@{}lcc@{}}
\toprule
\textbf{Rango RMSD} & \textbf{N} & \textbf{\%} \\
\midrule
$< \SI{0,5}{\angstrom}$ & 2.184 & 42,7\% \\
\SI{0,5}{\angstrom} -- \SI{1,0}{\angstrom} & 1.407 & 27,5\% \\
\SI{1,0}{\angstrom} -- \SI{2,0}{\angstrom} & 895 & 17,5\% \\
\SI{2,0}{\angstrom} -- \SI{5,0}{\angstrom} & 557 & 10,9\% \\
$> \SI{5,0}{\angstrom}$ & 72 & 1,4\% \\
\midrule
\textbf{Total} & \textbf{5.115} & \textbf{100\%} \\
\bottomrule
\end{tabular}
\end{table}

\subsection{RMSD por categor\'ia de tama\~no (mapeo SIFTS)}

\begin{table}[H]
\centering
\caption{RMSD desglosado por tama\~no proteico con mapeo SIFTS correcto.}
\label{tab:size_breakdown}
\begin{tabular}{@{}lrcccc@{}}
\toprule
\textbf{Categor\'ia} & \textbf{N} & \textbf{Media (\AA)} & \textbf{Mediana (\AA)} & \textbf{$<$ 1\AA} & \textbf{$<$ 2\AA} \\
\midrule
$<$ 50 res & 259 & 1,87 & 1,13 & 44,8\% & 66,4\% \\
50--100 & 1.156 & 0,94 & 0,57 & 72,2\% & 90,7\% \\
100--150 & 1.635 & 0,96 & 0,62 & 76,6\% & 90,9\% \\
150--200 & 151 & 2,30 & 1,18 & 36,4\% & 64,9\% \\
200--250 & 535 & 1,05 & 0,95 & 55,3\% & 91,6\% \\
250--300 & 662 & 1,12 & 0,74 & 63,3\% & 80,5\% \\
300--350 & 522 & 0,43 & 0,26 & 94,4\% & 99,2\% \\
350--400 & 192 & 1,70 & 0,61 & 64,6\% & 71,9\% \\
\bottomrule
\end{tabular}
\end{table}

\textit{Nota: Los datos corresponden al dataset validado (91 prote\'inas, 5.115 pares)
con mapeo SIFTS expl\'icito. Las variaciones entre categor\'ias reflejan la composici\'on
del dataset m\'as que una tendencia intr\'inseca del algoritmo. V\'ease la
Tabla~\ref{tab:rmsd_tamano} en la secci\'on de Resultados para los datos del
dataset diversificado de 10.432 prote\'inas \'unicas.}

\subsection{Impacto del mapeo SIFTS: comparaci\'on directa}

\begin{table}[H]
\centering
\caption{Comparaci\'on de m\'etricas con y sin mapeo SIFTS para los mismos 5.115 pares.}
\label{tab:sifts_comparison}
\begin{tabular}{@{}lcc@{}}
\toprule
\textbf{M\'etrica} & \textbf{Sin mapeo} & \textbf{Con mapeo SIFTS} \\
\midrule
RMSD medio & \SI{29,94}{\angstrom} & \SI{1,05}{\angstrom} \\
RMSD mediana & \SI{22,54}{\angstrom} & \SI{0,64}{\angstrom} \\
Desviaci\'on est\'andar & --- & \SI{1,37}{\angstrom} \\
$<$ \SI{0,5}{\angstrom} & 5,1\% & 42,7\% \\
$<$ \SI{1}{\angstrom} & 9,3\% & 70,2\% \\
$<$ \SI{2}{\angstrom} & 13,9\% & 87,7\% \\
$>$ \SI{5}{\angstrom} & 84,1\% & 1,4\% \\
$>$ \SI{10}{\angstrom} & 80,9\% & 0,4\% \\
\bottomrule
\end{tabular}
\end{table}
